\section{Pipeline de datos de Actor-Critic}

\subsection{Recolección, almacenamiento y distribución}

El flujo de datos del marco de Actor-Critic puede dividirse en cuatro etapas: recolección (rollout), almacenamiento (buffer), procesamiento (advantage estimation) y distribución (learning step). Cada etapa tiene sus trampas. Por ejemplo, si en la etapa de recolección todos los datos provienen de una sola deterministic policy, se produce low diversity, lo que provoca que en el replay buffer aparezcan muchas near-duplicate transitions, reduciendo el número efectivo de muestras de cada actualización real.

\begin{knowledgebox}{Coordinación entre rollout y buffer}
Para mantener la diversidad de la recolección, se suelen usar técnicas como epsilon/entropy noise o la mezcla de políticas aleatorias; al mismo tiempo, se procura evitar escribir la misma trayectoria de manera repetida en el buffer, haciendo al menos, al momento del enqueue, un tenure check (que registra cuántas veces se ha accedido a la transition).
\end{knowledgebox}

\subsection{Estimación de la ventaja y asignación de atribución}

La estimación de la ventaja (advantage estimation) es el centro del Actor-Critic. El algoritmo más usado es GAE($\lambda$), que hace un trade-off entre bias-variance:

$$\hat{A}_t = \sum_{l=0}^{T-t}(\gamma \lambda)^l \delta_{t+l},\quad \delta_t = r_t + \gamma V(s_{t+1}) - V(s_t)$$

Elegir $\lambda=0.95$ es un valor por defecto común, y al mismo tiempo es necesario mantener estable el entrenamiento de la función de valor $\hat{V}$; de lo contrario, la varianza del advantage se dispara.

\subsection{Resumen de la sección}

El pipeline de datos de Actor-Critic abarca rollout, buffer, advantage estimation y update step. Cada eslabón necesita monitoreo para evitar la instability provocada por desviaciones abruptas.

